\scanpage{160}
La solution du pb \inserted{\unsure{ci-dessus|semi-linéaire}$^{*}$} universel est un module quotient de $\bigotimes E_i^0\oplus\bigotimes E_i$, c'est une extension de $\bigotimes E_i/E_i^0$ par $\bigotimes E_i^0$ :
\[
0\longrightarrow\bigotimes E_i^0\xrightarrow{i}\mathop{*}_{i}(E_i,\sigma_i)^{\deleted{\unsure{0|\wedge}}}\xrightarrow{j}\bigotimes(E_i/E_i^0)\longrightarrow0
\]
% slanted
\margin{attention, ce n'est peut-être vrai tel \unsure{quel|que} si on suppose p. ex. \unsure{que} $E_i/E_i^0$ \unsure{et} $E_i^0$ \unsure{plat. des|plats. des} extensions convenablement \unsure{scindées|scindée} et $2$ inv.}
% slanted
\margin{NB On devrait écrire $E_i^-$ al d $E_i^0$ et $E_i^+$ al d $E_i/E_i^0$}

Quand on a des isos $E_i/E_i^0\simeq k$ alors $\mathop{*}_{i}(E_i,\sigma_i)$ apparait comme une extension de \deleted{$k$}$k$ par $\bigotimes E_i^0$ -- on retrouve le contexte des espaces affines.
% slanted
\margin{$*$ \deleted{En effet,} Car on a si \deleted{deux parmi les $\sigma_i$ sont $\sigma_i$} \deleted{sont} $\sigma_i=\mathrm{id}$, alors \deleted{$j$ est iso, alors} \deleted{\unsure{nul}!} $\mathop{\unsure{*|\bigotimes}}_{i}(E_i,\sigma_i)\simeq\bigotimes E_i/E_i^0\oplus\bigotimes E_i^0$ et $i$ est quand même injectif}

On a toujours un hom canonique
\[
\mathop{*}_{i}(E_i,\sigma_i)\xrightarrow{\tau}\bigotimes E_i
\]
donné par la formule devinée tantôt (pour $n=2$, c'est la formule
\[
x_1*x_2=\sigma_1 x_1\otimes x_2+x_1\otimes\sigma_2 x_2\qquad)
\]
mais on a
\diagram{\begin{tikzcd}[ampersand replacement=\&,column sep=large]
\bigotimes E_i^0 \arrow[d,"i"'] \arrow[dr,"{(-2)^{n-1}\mathrm{can}}"] \& \\
\mathop{*}(E_i,\sigma_i) \arrow[r,"\tau"] \arrow[dr,"{2^{n-1}j}"'] \& \bigotimes E_i \arrow[d,"\mathrm{can}"]\\
\& \bigotimes E_i/E_i^0
\end{tikzcd}}
% slanted
\margin{finalement, cette application $\tau$ n'est pas tellement utile, je crois, si $2$ n'est pas supposé inversible.}
C'est donc inj. si $2$ inv., mais pas en général.

On a là une loi de composition \inserted{interne} associative, commutative,
% archive sheet number: 159

